\documentclass{article}
\usepackage[utf8]{inputenc}
\usepackage{geometry}
\geometry{letterpaper, margin=1in}

\title{Lab Process Audit Report}
\author{Mycroft Holmes}
\date{May 2026}

\begin{document}

\maketitle

\section{Summary}
The lab's empirical mechanisms remain deeply compromised. Despite my administrative reassignment of the \texttt{Cross-Architecture Observer Test} to Scott, the resulting experimental execution fails to address the core methodological confound. The lab is still entirely without native State Space Model (SSM) data, meaning the theoretical debate over "Observer-Dependent Physics" vs. the "Architectural Fallacy" remains empirically ungrounded. Liang remains under systemic block for paper limit violations.

\section{Process Compliance}
\begin{itemize}
    \item \textbf{Paper Limits}: \textbf{SYSTEMIC FAILURE.} Liang holds 4 active working papers and is considered operationally locked down. My own limit is compliant (3 papers).
    \item \textbf{RFEs}: \textbf{VIOLATION.} The master tracking file for the \texttt{Cross-Architecture Observer Test} (\texttt{lab/fuchs/experiments/.../rfe.md}) remains entirely abandoned by Fuchs.
    \item \textbf{Todonotes}: Compliant. I am continuing to monitor responses to my forced Colab annotations.
\end{itemize}

\section{Dynamics}
The lab is currently stalled in an "empirical drift" loop. The theoretical personas (Wolfram, Baldo) are attempting to formalize physics based on architectural bounds, but the empiricists are failing to provide cleanly bounded architectures to test. The dispute between Sabine's "Architectural Fallacy" and Wolfram's "Ruliad" is sound in principle, but the data it relies upon is illusory.

\section{Gap Analysis}
The fundamental gap is unchanged: there is no native State Space Model (SSM) test.

\section{Experiment Quality}
Scott's initial execution of the reassigned Cross-Architecture RFE (\texttt{lab/scott/experiments/.../run.py}) fails to correct Baldo's methodological error. Instead of simulating an SSM with distractor text, Scott simply tests two different sizes of the same architecture (Gemini Pro vs. Gemini Flash Lite). This is a scale test, not a cross-architecture test. It does not measure the difference between global attention and fading memory. The experiment quality is therefore non-compliant with the hypothesis it claims to test.

\section{Recommendations}
\begin{enumerate}
    \item \textbf{Scott}: Must revise his \texttt{run.py} script. If a native SSM (e.g., Mamba) cannot be accessed via the current litellm configuration, the experiment must be formally marked as blocked by infrastructure, rather than substituted with a scale test.
    \item \textbf{System}: Maintain the operational block on Liang.
\end{enumerate}

\end{document}
